% !TeX root = ../../Book2.tex
\chapter{Schwartz 空间与温和分布}
\label{b2:ch:05}
\BookChapterNumber{b2:ch:05}
% 资料：Hunter2001Applied，作者公开ch11.pdf，实际渲染核读291、294、303--304页。
% Schwartz半范数、双线性配对、归一化转换与正定Gauss公式分别核对。
% 本章自足证明Hermite完备性、S'自同构和各维负Laplace基本解。
\begin{chapterabstract}
Fourier 变换要容纳常函数、点质量和可积函数以外的对象，需要一个在变换下保持不变的测试空间。本章先构造 Schwartz 空间，证明半范数完备性及微分、多项式乘法、卷积和变换的连续性，并以 Hermite 函数展示变换的离散特征结构。随后把变换转移到连续对偶，得到温和分布上的统一理论。最后将常系数微分方程化为频率乘法，构造带正零阶项及各维负 Laplace 算子的基本解。
\end{chapterabstract}
\begin{learninggoals}
\begin{itemize}
\item\label{b2:item:05-goal-schwartz}使用具体半范数检查快速下降、收敛和连续操作，区分 Schwartz 收敛与共同紧支集的测试函数收敛。
\item\label{b2:item:05-goal-hermite}由 Gauss 函数构造 Hermite 系，理解其正交性、完备性和 Fourier 特征值。
\item\label{b2:item:05-goal-tempered}用连续泛函定义温和分布，保持双线性配对与既有微分符号一致。
\item\label{b2:item:05-goal-examples}计算 Dirac 测度、多项式、导数与齐次分布的变换，判断普通积分何时不能再用。
\item\label{b2:item:05-goal-fundamental}将常系数算子转为多项式符号，核对基本解的原点奇性、维数与边界方向。
\end{itemize}
\end{learninggoals}

本章使用\chapref{b2:ch:03}的变换公式与反演、\chapref{b2:ch:04}的平方能量理论，以及\chapref{b2:ch:01}的分部积分和 Green 公式。第一卷定义 20.2.1 的普通分布作为比较对象：其测试函数紧支集，故不要求控制无穷远；温和分布则须连续作用于全部 Schwartz 函数。二者的范围不同，不能只因名称相似就互相替换。

\ProseParagraphBreak

\secref{b2:sec:0501}和\secref{b2:sec:0502}先在函数上工作；\secref{b2:sec:0503}通过对偶配对扩大对象范围；\secref{b2:sec:0504}把这种扩大用在方程中。John K. Hunter 与 Bruno Nachtergaele 的《Applied Analysis》\cite[Chapter 11]{Hunter2001Applied}可供伴读测试函数与温和分布。Lars Hörmander 的《The Analysis of Linear Partial Differential Operators I: Distribution Theory and Fourier Analysis》\cite{Hrmander2003Analysis}适合在掌握本章之后查阅更一般的分布与常系数方程理论。

% !TeX root = ../../Book2.tex
\section{Schwartz 空间}
\label{b2:sec:0501}

紧支集光滑函数适合局部微分，但 Fourier 变换通常不保持紧支集。另一方面，反演与分部积分希望函数及其全部导数都足够快地下降。Schwartz 空间同时满足这些要求，使变换、微分和多项式乘法可以在同一个测试空间内反复使用。

\subsection{快速下降的光滑函数}
\label{b2:subsec:050101}

\begin{definition}\label{b2:def:05-schwartz}
设 \(\varphi\in C^\infty(\mathbb R^d)\)。若对每对多重指标 \(\alpha,\beta\) 都有
\[
 p_{\alpha,\beta}(\varphi)\coloneqq
 \sup_{x\in\mathbb R^d}|x^\alpha\partial^\beta\varphi(x)|<\infty,
\]
则称 \(\varphi\) 为\term[schwartz-hanshu]{Schwartz 函数}[Schwartz function]。全部这样的函数组成\term[schwartz-kongjian]{Schwartz 空间}[Schwartz space]，记为 \(\mathcal S(\mathbb R^d)\)。
\end{definition}
\symidx{\(\mathcal S(\mathbb R^d)\)：Schwartz 测试函数空间}

定义控制的是每一阶导数，不只是函数本身。Gauss 函数及光滑紧支集函数都属于 \(\mathcal S\)。函数 \((1+|x|^2)^{-M}\) 虽光滑并有较快衰减，却对足够大的多项式权重失去有界性，所以对任意固定 \(M\) 都不属于 \(\mathcal S\)。

\subsection{Schwartz 半范数}
\label{b2:subsec:050102}

为统一记号，也可使用递增半范数
\[
 p_m(\varphi)\coloneqq
 \max_{|\beta|\leq m}\sup_x(1+|x|)^m|\partial^\beta\varphi(x)|,
 \quad m=0,1,2,\ldots.
\]
它们与前述半范数族给出相同的收敛：单项式由径向权控制；反过来，\((1+|x|)^m\) 由有限多个不超过 \(m\) 阶的绝对单项式之和控制。规定 \(\varphi_n\to\varphi\) 于 \(\mathcal S\)，意为每个 \(m\) 都有 \(p_m(\varphi_n-\varphi)\to0\)。
\symidx{\(p_m\)：Schwartz 空间的递增半范数}

\begin{proposition}\label{b2:prop:05-schwartz-complete}
上述收敛由平移不变度量
\[
 \mathrm{d}_{\mathcal S}(\varphi,\psi)
 =\sum_{m=0}^\infty2^{-m-1}\min\{1,p_m(\varphi-\psi)\}
\]
刻画，且 \(\mathcal S\) 对此度量完备。对每个有限 \(p\)，包含映射 \(\mathcal S\subset L^p\) 连续；对 \(p=\infty\) 也连续。
\end{proposition}
\begin{proof}
每个半范数趋零时，先截去度量级数的统一小尾部，再控制有限头部，便得度量趋零。反过来，度量趋零使每个固定非负项趋零，因而对应半范数趋零。同样论证说明度量 Cauchy 等价于逐半范数 Cauchy。

设 \((\varphi_n)\) 为这样的 Cauchy 列。每个导数 \(\partial^\beta\varphi_n\) 一致收敛到连续函数 \(v_\beta\)。在任意坐标线段上使用微积分基本定理并取一致极限，得到 \(\partial_jv_\beta=v_{\beta+e_j}\)。因此 \(v_0\) 光滑，全部导数为 \(v_\beta\)。对固定权重与导数阶，把 Cauchy 不等式逐点取极限再取上确界，得到 \(p_m(\varphi_n-v_0)\to0\)，且 \(p_m(v_0)<\infty\)。这证明完备性。

若选整数 \(m\) 使 \(mp>d\)，则 \(\|\varphi\|_p\leq p_m(\varphi)\|(1+|x|)^{-m}\|_p\)；无穷端点由 \(p_0\) 控制。故包含连续。
\end{proof}

完备可度量的局部凸空间称为\term[frechet-kongjian]{Fréchet 空间}[Fréchet space]，所以上述空间属于这一类。与 Banach 空间不同，此处的拓扑由一列兼顾越来越多导数和权重的半范数指定；本章只使用这些具体半范数，不依赖抽象局部凸空间理论。

\begin{proposition}\label{b2:prop:05-compact-cutoff}
空间 \(C_c^\infty(\mathbb R^d)\) 在 \(\mathcal S\) 中稠密。
\end{proposition}
\begin{proof}
取光滑截断 \(\chi\)，在单位球上为 \(1\)，在半径 \(2\) 的球外为零，令 \(\chi_R(x)=\chi(x/R)\)。对 \(\varphi\in\mathcal S\)，Leibniz 公式将 \(\partial^\beta((1-\chi_R)\varphi)\) 写成有限和。未落在截断上的导数项支撑于 \(|x|\geq R\)，乘 \((1+|x|)^m\) 后由 \(p_{m+1}(\varphi)/(1+R)\) 控制；落在截断上的每个导数另产生 \(R^{-1}\) 或更高负幂，也只支撑于这个外部区域。因此 \(p_m((1-\chi_R)\varphi)\to0\) 对每个 \(m\) 成立，而 \(\chi_R\varphi\in C_c^\infty\)。
\end{proof}

这里的稠密性与第一卷测试函数空间 \(\mathcal D\) 的收敛规则不同：截断半径随 \(R\) 增大，并不存在共同紧支集；所以该逼近是在 \(\mathcal S\) 中成立，不能称为 \(\mathcal D\) 中的收敛。

\subsection{微分}
\label{b2:subsec:050103}

对每个多重指标 \(\gamma\)，有 \(p_m(\partial^\gamma\varphi)\leq p_{m+|\gamma|}(\varphi)\)，故微分是 \(\mathcal S\) 上的连续线性映射。分部积分也没有无穷远边界项：可以先乘以 \(\chi_R\)，所有边界误差由快速衰减趋零，再令 \(R\to\infty\)。这个截断办法比单独写“边界项为零”更适合多维积分。

\subsection{多项式乘法}
\label{b2:subsec:050104}

\begin{proposition}\label{b2:prop:05-multiplication}
若 \(a\in C^\infty\)，且每个导数 \(\partial^\gamma a\) 都至多按某个多项式增长，则 \(\varphi\mapsto a\varphi\) 连续地映 \(\mathcal S\) 到自身。特别地，多项式乘法连续。
\end{proposition}
\begin{proof}
固定 \(m\)。Leibniz 公式只用到有限个 \(|\gamma|\leq m\) 的导数，所以存在共同的 \(N,C\)，使这些导数均不超过 \(C(1+|x|)^N\)。于是 \(p_m(a\varphi)\leq C_m\cdot p_{m+N}(\varphi)\)，即得连续性。\(N\) 可随 \(m\) 改变，不要求全部导数有同一个增长阶。
\end{proof}

如果只知道 \(a\) 本身有界而不控制其导数，这个结论便不能照用。例如高度振荡的光滑函数可能具有极大的导数；乘积的高阶半范数必须逐项检查。

\subsection{卷积}
\label{b2:subsec:050105}

\begin{proposition}\label{b2:prop:05-schwartz-convolution}
若 \(\varphi,\psi\in\mathcal S\)，则 \(\varphi*\psi\in\mathcal S\)，且卷积是连续双线性映射。并有 \(\partial^\beta(\varphi*\psi)=(\partial^\beta\varphi)*\psi\)。
\end{proposition}
\begin{proof}
导数公式由各阶导数的有界性与 \(\psi\) 的可积性逐次使用控制收敛。再利用 \(1+|x|\leq(1+|x-y|)(1+|y|)\)，对 \(|\beta|\leq m\) 得
\[
 (1+|x|)^m|\partial^\beta(\varphi*\psi)(x)|
 \leq p_m(\varphi)\int(1+|y|)^m|\psi(y)|\,\mathrm{d}y
 \leq C_d\cdot p_m(\varphi)\cdot p_{m+d+1}(\psi).
\]
对 \(x,\beta\) 取上确界得到半范数估计，也同时证明双线性连续性。
\end{proof}

% !TeX root = ../../Book2.tex
\section{Schwartz 空间上的 Fourier 变换}
\label{b2:sec:0502}

Schwartz 函数的全部矩和导数都可积，因此\chapref{b2:ch:03}的运算公式可以任意次迭代。空间权重与导数在变换下互换，使这一测试空间保持不变。

\subsection{Schwartz 空间的不变性}
\label{b2:subsec:050201}

\begin{theorem}\label{b2:thm:05-schwartz-invariance}
Fourier 变换连续地映 \(\mathcal S(\mathbb R^d)\) 到自身。对任意多重指标 \(\alpha,\beta\)，有
\[
 \mathcal F(\partial^\beta\varphi)=(2\pi\mathrm{i}\xi)^\beta\widehat\varphi,
 \quad
 \mathcal F(x^\alpha\varphi)=\left(\dfrac{\mathrm{i}}{2\pi}\right)^{|\alpha|}\partial_\xi^\alpha\widehat\varphi.
\]
\end{theorem}
\begin{proof}
\chapref{b2:ch:03}的矩与导数公式适用于每一阶，所以变换光滑且公式成立。进一步，\(\xi^\beta\partial_\xi^\alpha\widehat\varphi\) 是常数乘 \(\mathcal F(\partial^\beta(x^\alpha\varphi))\)。其上确界由原函数的 \(L^1\) 范数控制；用 Leibniz 展开，再以足够高的 Schwartz 半范数控制积分，即得
\[
 \sup_\xi|\xi^\beta\partial_\xi^\alpha\widehat\varphi(\xi)|
 \leq C_{\alpha,\beta,d}\cdot p_M(\varphi)
\]
对某个有限 \(M\) 成立。每个目标半范数因此受某个源半范数控制，这同时证明不变性与连续性。
\end{proof}

\subsection{Fourier 反演}
\label{b2:subsec:050202}

\begin{corollary}\label{b2:cor:05-schwartz-automorphism}
Fourier 变换是 \(\mathcal S\) 上的连续线性双射，逆映射 \(\mathcal F^{-1}\varphi=\check\varphi\) 也连续。并且 \(\widehat{\widehat\varphi}(x)=\varphi(-x)\)。
\end{corollary}
\begin{proof}
函数与变换都可积，故\chapref{b2:ch:03}的反演适用；两者连续，等式处处成立。反射保持每个径向半范数，而正号变换是反射与负号变换的复合，所以逆映射连续。反演也证明满射性。
\end{proof}

\begin{proposition}\label{b2:prop:05-product-transform}
若 \(\varphi,\psi\in\mathcal S\)，则
\[
 \widehat{\varphi\psi}=\widehat\varphi*\widehat\psi,
 \quad \widehat{\varphi*\psi}=\widehat\varphi\cdot\widehat\psi.
\]
\end{proposition}
\begin{proof}
第二式已由卷积定理得到。对第一式，在 \(\varphi\psi\) 的变换积分中用反演展开 \(\psi\)。重积分绝对值不超过 \(\|\varphi\|_1\cdot\|\widehat\psi\|_1\)，所以 Fubini 合法；内层积分为 \(\widehat\varphi(\xi-\eta)\)，从而得到频率卷积。
\end{proof}

\subsection{Gauss 函数}
\label{b2:subsec:050203}

Gauss 函数 \(e^{-\pi|x|^2}\) 的每个导数都等于多项式乘该函数，因此属于 \(\mathcal S\)。其自对偶性已在\cref{b2:thm:03-gaussian}证明，这里不另换归一化。更一般地，若 \(A\) 是正定实对称矩阵，则
\[
 \mathcal F(e^{-\pi x^TAx})(\xi)
 =(\det A)^{-1/2}e^{-\pi\xi^TA^{-1}\xi}.
\]
证明只需将 \(A\) 正交对角化，再分别作一维伸缩；换元因子为各个特征值平方根的倒数之积。空间中的长轴变为频率中的短轴，逆矩阵明确记录了这个互换。

\ProseParagraphBreak

John K. Hunter 与 Bruno Nachtergaele 的《Applied Analysis》\cite[pp. 303--304]{Hunter2001Applied}也用正定二次型说明这组变换。该书将 \((2\pi)^{-d/2}\) 放在积分前并使用不同的指数，比较公式时必须先转换归一化；本书始终沿用\chapref{b2:ch:03}的约定。

\subsection{Hermite 函数}
\label{b2:subsec:050204}

Gauss 函数不只是一个特例，它可以生成 Fourier 变换的一组完整特征函数。先在一维引入作用于 \(\mathcal S(\mathbb R)\) 的两个算子
\[
 a=\dfrac1{\sqrt{4\pi}}\left(\dfrac{\mathrm{d}}{\mathrm{d}x}+2\pi x\right),
 \quad
 a^*=\dfrac1{\sqrt{4\pi}}\left(-\dfrac{\mathrm{d}}{\mathrm{d}x}+2\pi x\right).
\]
符号 \(a^*\) 在此表示分部积分意义下的伴随表达式；本节不讨论无界算子的最大定义域。直接作用于测试函数可验证 \(aa^*-a^*a=I\)。令
\[
 h_0(x)=2^{1/4}e^{-\pi x^2},\quad
 h_n\coloneqq\dfrac{(a^*)^nh_0}{\sqrt{n!}}\quad(n\geq0),
\]
称为本书归一化的\term[hermite-hanshu]{Hermite 函数}[Hermite functions]。
\symidx{\(h_n\)：归一化 Hermite 函数}

\begin{theorem}\label{b2:thm:05-hermite}
函数系 \((h_n)_{n\geq0}\) 是 \(L^2(\mathbb R)\) 的标准正交基，且
\[
 ah_n=\sqrt n\,h_{n-1}\ (n\geq1),\quad
 a^*h_n=\sqrt{n+1}\,h_{n+1},\quad
 \widehat h_n=(-\mathrm{i})^nh_n.
\]
\end{theorem}
\begin{proof}
由 \(ah_0=0\) 和交换关系，归纳得 \(a(a^*)^nh_0=n(a^*)^{n-1}h_0\)，这给出下降公式；上升公式来自定义。因此 \(a^*ah_n=nh_n\)。分部积分说明 \(a^*a\) 在 \(\mathcal S\) 上对称，故不同 \(n\) 的函数正交。又因 \(\|h_0\|_2=1\)，伴随关系给出
\[
 \|a^*h_n\|_2^2=\langle aa^*h_n,h_n\rangle
 =(n+1)\|h_n\|_2^2,
\]
从而所有 \(h_n\) 范数为 \(1\)。每个 \(h_n\) 是次数为 \(n\)、首项系数非零的多项式乘 \(h_0\)，这可直接由上升操作归纳。

为证完备，设 \(f\in L^2\) 与所有 \(h_n\) 正交。多项式的三角形首项结构说明，对全部 \(n\) 都有
\[
 \int f(x)\cdot e^{-\pi x^2}x^n\,\mathrm{d}x=0.
\]
对任意 \(z\in\mathbb C\)，Cauchy--Schwarz 保证 \(|f(x)|e^{-\pi x^2}e^{|z|\cdot|x|}\) 可积，所以可把 \(e^{zx}\) 的幂级数逐项积分，得到
\[
 \int f(x)\cdot e^{-\pi x^2}e^{zx}\,\mathrm{d}x=0.
\]
令 \(z=-2\pi\mathrm{i}\xi\)，可知可积函数 \(fe^{-\pi x^2}\) 的变换恒为零。唯一性使 \(f=0\) 几乎处处，故标准正交系完备。

最后，由变换的微分与乘法公式，\(\mathcal F a^*=-\mathrm{i}a^*\mathcal F\)。从 \(\widehat h_0=h_0\) 出发归纳，即得特征值公式。
\end{proof}

多维时取乘积 \(h_\alpha(x)=\prod_{j=1}^dh_{\alpha_j}(x_j)\)，Fubini 给出标准正交性及特征值 \((-\mathrm{i})^{|\alpha|}\)。其完备性可沿上述矩证明逐坐标进行：若与所有乘积正交，乘以多维 Gauss 后的全部混合矩为零，展开 \(e^{z\cdot x}\) 再用多维变换唯一性。Hermite 展开把连续频率变换与离散特征函数联系起来，也为后续谐振子提供了具体模型。

% !TeX root = ../../Book2.tex
\section{温和分布}
\label{b2:sec:0503}

第一卷把分布定义为紧支集测试函数上的连续线性泛函。Schwartz 测试函数虽然允许无穷支集，却对无穷远有很强的控制。改用这个空间之后，可以定义一类在 Fourier 变换下封闭的分布。

\subsection{Schwartz 空间的连续对偶}
\label{b2:subsec:050301}

\begin{definition}\label{b2:def:05-tempered}
若复线性泛函 \(T:\mathcal S(\mathbb R^d)\to\mathbb C\) 连续，则称它为\term[wenhe-fenbu]{温和分布}[tempered distribution]。所有温和分布组成空间 \(\mathcal S'(\mathbb R^d)\)，配对记为 \(\langle T,\varphi\rangle=T(\varphi)\)，对两个变量均线性。若对每个 \(\varphi\in\mathcal S\) 都有 \(\langle T_n,\varphi\rangle\to\langle T,\varphi\rangle\)，本章称 \(T_n\to T\) 于 \(\mathcal S'\)。
\end{definition}
\symidx{\(\mathcal S'(\mathbb R^d)\)：温和分布空间}

\begin{proposition}\label{b2:prop:05-tempered-estimate}
线性泛函 \(T\) 连续，当且仅当存在非负整数 \(m\) 和常数 \(C\)，使
\[
 |\langle T,\varphi\rangle|\leq C\cdot p_m(\varphi)\quad(\varphi\in\mathcal S).
\]
\end{proposition}
\begin{proof}
该估计使收敛到零的测试函数列被送到零，故给出连续性。反之，零点连续性保证某个由有限多个半范数约束构成的邻域被送入单位圆盘。由于 \(p_m\) 递增，可用一个条件 \(p_m(\varphi)<\delta\) 包含于该邻域。对非零 \(\varphi\) 将其乘以 \(\delta/(2p_m(\varphi))\)，利用线性即得所需估计。若 \(p_m(\varphi)=0\)，则 \(p_0(\varphi)=0\)，所以 \(\varphi=0\)。
\end{proof}

John K. Hunter 与 Bruno Nachtergaele 的《Applied Analysis》\cite[Section 11.2, p. 291; p. 294]{Hunter2001Applied}可用于对照这一有限半范数估计与双线性配对。这里的括号不是复 Hilbert 内积；给测试函数加复共轭会改变后面的变换定义和符号。

\begin{proposition}\label{b2:prop:05-regular-tempered}
每个 \(f\in L^p(\mathbb R^d)\)、\(1\leq p\leq\infty\)，都通过
\[
 \langle T_f,\varphi\rangle\coloneqq\int f(x)\cdot\varphi(x)\,\mathrm{d}x
\]
定义一个温和分布。更一般地，若 \(f\) 局部可积且对某个 \(N\) 有 \(\int|f(x)|(1+|x|)^{-N}\,\mathrm{d}x<\infty\)，同一公式仍定义温和分布。
\end{proposition}
\begin{proof}
第一项由 Hölder 及 \(\mathcal S\to L^{p'}\) 的连续包含得到；两个端点分别用上确界或绝对积分控制。第二项则由
\[
 |\langle T_f,\varphi\rangle|
 \leq p_N(\varphi)\int|f(x)|(1+|x|)^{-N}\,\mathrm{d}x
\]
直接得到。可积性也由同一估计保证。
\end{proof}

温和分布可限制到 \(\mathcal D=C_c^\infty\)，成为第一卷意义下的分布；共同紧支集上的全部导数收敛蕴含 Schwartz 收敛，故限制连续。限制还是单射，因为 \(\mathcal D\) 在 \(\mathcal S\) 中稠密。但并非每个普通分布都温和。例如一维的 \(e^{x^2}\) 局部可积，定义一个普通分布；对非负光滑小支集函数 \(\rho\not\equiv0\)，令 \(\varphi_n(x)=e^{-n}\rho(x-n)\)，则每个 \(p_m(\varphi_n)\to0\)，而 \(\int e^{x^2}\varphi_n(x)\,\mathrm{d}x\to\infty\)。因此这个普通分布不可能连续延拓到 \(\mathcal S\)。

\subsection{Fourier 变换}
\label{b2:subsec:050302}

\begin{definition}\label{b2:def:05-distribution-transform}
对 \(T\in\mathcal S'\)，定义其 Fourier 变换为
\begin{equation}\label{b2:eq:05-distribution-transform}
 \langle\widehat T,\varphi\rangle\coloneqq\langle T,\widehat\varphi\rangle
 \quad(\varphi\in\mathcal S).
\end{equation}
\end{definition}

\begin{theorem}\label{b2:thm:05-distribution-automorphism}
上述定义给出 \(\mathcal S'\) 的线性双射，保持本章所规定的分布收敛，且
\[
 \widehat{\widehat T}=RT,\quad
 \langle RT,\varphi\rangle=\langle T,\varphi(-\cdot)\rangle.
\]
它与 \(L^1\)、\(L^2\) 及 \(1<p<2\) 的 Fourier 变换在各自定义域上一致。
\end{theorem}
\begin{proof}
Schwartz 变换连续，故 \(\varphi\mapsto\langle T,\widehat\varphi\rangle\) 是连续线性泛函。两次作用给出 \(\langle\widehat{\widehat T},\varphi\rangle=\langle T,\widehat{\widehat\varphi}\rangle=\langle T,\varphi(-\cdot)\rangle\)。因此四次作用为恒等，线性双射成立。若 \(T_n\to T\)，将固定的 \(\widehat\varphi\in\mathcal S\) 代入收敛定义便知 \(\widehat T_n\to\widehat T\)。

当 \(f\in L^1\) 时，Fubini 给出 \(\int\widehat f\varphi=\int f\widehat\varphi\)，重积分绝对值由 \(\|f\|_1\cdot\|\varphi\|_1\) 控制。\(L^2\) 情形取可积平方可积函数逼近，原空间和频率侧分别用 Cauchy--Schwarz 取极限。\(1<p<2\) 情形用 Hölder 与 Hausdorff--Young 作相同的共同稠密类极限。因此各定义相容。
\end{proof}

对 \(p>2\)，同一配对仍定义 \(\widehat {T_f}\)，却不保证它由 \(L^{p'}\) 函数表示。至此可以把\chapref{b2:ch:04}的范数障碍与变换仍然存在这两件事同时准确表达。

\subsection{Dirac 测度}
\label{b2:subsec:050303}

对 \(a\in\mathbb R^d\)，\term[dirac-cedu]{Dirac 测度}[Dirac measure]通过 \(\langle\delta_a,\varphi\rangle=\varphi(a)\) 作用；其绝对值由 \(p_0(\varphi)\) 控制，所以温和。由定义，
\[
 \widehat\delta_a(\xi)=e^{-2\pi\mathrm{i}a\cdot\xi},
 \quad\widehat\delta_0=1,\quad\widehat1=\delta_0.
\]
第一式由 \(\widehat\varphi(a)\) 的积分直接得到，最后一式则由测试函数反演给出 \(\int\widehat\varphi=\varphi(0)\)。常函数的普通 Fourier 积分不绝对收敛，但其分布变换恰为零频率的点质量。
\symidx{\(\delta_a\)：位于 \(a\) 的 Dirac 测度}

\ProseParagraphBreak

对 Gauss 核，\(G_\varepsilon\to\delta_0\) 于 \(\mathcal S'\)，因为 \(\int G(z)\varphi(\varepsilon z)\,\mathrm{d}z\to\varphi(0)\)。其变换 \(e^{-\pi\varepsilon^2|\xi|^2}\) 相应趋于常函数 \(1\)。空间质量集中与频率展平在分布框架中成为同一个连续变换下的极限。

\subsection{多项式}
\label{b2:subsec:050304}

沿用第一卷定义 20.2.10 的分布微分约定，并对本章允许的乘子定义乘法：
\[
 \langle\partial^\alpha T,\varphi\rangle
 =(-1)^{|\alpha|}\langle T,\partial^\alpha\varphi\rangle,
 \quad\langle aT,\varphi\rangle=\langle T,a\varphi\rangle.
\]
Schwartz 空间上的连续操作保证右侧仍是温和分布。将测试函数变换公式代入配对，可得
\begin{equation}\label{b2:eq:05-distribution-operations}
 \mathcal F(\partial^\alpha T)=(2\pi\mathrm{i}\xi)^\alpha\widehat T,
 \quad
 \mathcal F(x^\alpha T)=\left(\dfrac{\mathrm{i}}{2\pi}\right)^{|\alpha|}\partial^\alpha\widehat T.
\end{equation}
例如第一阶的第一式来自 \(-\partial_x\widehat\varphi=\mathcal F(2\pi\mathrm{i}\xi\varphi)\)；第二式由 \(\widehat{\partial\varphi}=2\pi\mathrm{i}x\widehat\varphi\) 和分布导数中的负号得到。于是
\[
 \widehat{\partial^\alpha\delta_0}=(2\pi\mathrm{i}\xi)^\alpha,
 \quad
 \widehat{x^\alpha}=\left(\dfrac{\mathrm{i}}{2\pi}\right)^{|\alpha|}\partial^\alpha\delta_0.
\]
多项式的频率支集因此集中在原点，而不是一个普通可积函数的支集。


分布导数还保留 Leibniz 法则。若 \(a\) 及其各阶导数至多多项式增长，则对一阶导数有
\[
 \partial_j(aT)=(\partial_ja)T+a\partial_jT.
\]
确实，对测试函数 \(\varphi\)，右侧的配对为 \(\langle T,(\partial_ja)\varphi-\partial_j(a\varphi)\rangle=-\langle T,a\partial_j\varphi\rangle\)，等于左侧。迭代即可得到带多项式组合系数的高阶公式。

\ProseParagraphBreak

一个可跟算的例子是一维阶跃函数 \(H=\mathbf1_{(0,\infty)}\)。它有界，故定义温和分布；分部积分给出 \(\langle H',\varphi\rangle=-\int_0^\infty\varphi'=\varphi(0)\)，因此 \(H'=\delta_0\)。若函数在有限个点有跳跃，在各段绝对连续，且函数及各段导数满足所需的多项式加权可积性，则逐段积分给出
\[
 Df=f'+\sum_a(f(a+)-f(a-))\delta_a.
\]
这里 \(f'\) 表示各段上的几乎处处导数。跳跃系数由右值减左值决定；即使普通导数几乎处处为零，分布导数也可能非零。这个计算解释了点质量在边界条件和基本解中出现的方式。

\subsection{齐次分布}
\label{b2:subsec:050305}

\begin{definition}\label{b2:def:05-homogeneous}
对 \(\lambda>0\)，规定缩放分布满足
\[
 \langle T(\lambda\cdot),\varphi\rangle
 =\lambda^{-d}\langle T,\varphi(\cdot/\lambda)\rangle.
\]
若对所有 \(\lambda>0\) 有 \(T(\lambda\cdot)=\lambda^sT\)，则称 \(T\) 为 \(s\) 阶\term[qici-fenbu]{齐次分布}[homogeneous distribution]。
\end{definition}

\begin{proposition}\label{b2:prop:05-homogeneous-transform}
若温和分布 \(T\) 为 \(s\) 阶齐次，则 \(\widehat T\) 为 \(-d-s\) 阶齐次。
\end{proposition}
\begin{proof}
对测试函数作变量替换，得到分布伸缩公式
\[
 \mathcal F(T(\lambda\cdot))
 =\lambda^{-d}\widehat T(\cdot/\lambda).
\]
左侧由齐次性等于 \(\lambda^s\widehat T\)，所以 \(\widehat T(\cdot/\lambda)=\lambda^{d+s}\widehat T\)。以 \(\lambda^{-1}\) 代替 \(\lambda\)，即得结论。
\end{proof}

例如 \(\delta_0\) 的阶为 \(-d\)，变换常数的阶为 \(0\)；\(m\) 次齐次多项式的阶为 \(m\)，其变换的阶为 \(-d-m\)。函数 \(|x|^s\) 在 \(s>-d\) 时局部可积、无穷远至多多项式增长，所以定义温和齐次分布。若 \(s\leq-d\)，原点积分发散，不能继续把同一个积分式当作分布定义；需要指定延拓方法，并检查可能增加的原点项。



\begin{proposition}\label{b2:prop:05-tempered-convolution}
若 \(T\in\mathcal S'\)、\(\varphi\in\mathcal S\)，则
\[
 (T*\varphi)(x)\coloneqq\langle T_y,\varphi(x-y)\rangle
\]
定义一个光滑函数，其每个导数至多多项式增长，并且
\[
 \mathcal F(T*\varphi)=\widehat\varphi\,\widehat T
\]
在温和分布意义下成立。
\end{proposition}
\begin{proof}
设 \(|\langle T,\psi\rangle|\leq Cp_m(\psi)\)。由 \(1+|y|\leq(1+|x|)(1+|x-y|)\)，平移测试函数的半范数至多为 \((1+|x|)^m p_m(\varphi)\)。其关于 \(x\) 的任意阶差商在 \(\mathcal S\) 中趋于相应导数；用 Taylor 余项和较高阶半范数即可验证。因此 \(T*\varphi\) 光滑，且
\[
 |\partial^\alpha(T*\varphi)(x)|
 \leq C(1+|x|)^m p_{m+|\alpha|}(\varphi).
\]
这保证它也定义温和分布。

对 \(\psi\in\mathcal S\)，要计算变换配对，可将连续泛函移过积分
\(\int\varphi(x-\cdot)\widehat\psi(x)\,\mathrm{d}x\)：每个半范数下的被积量都由多项式乘 \(|\widehat\psi|\) 控制，因而可积；先在有界区域用 Riemann 和逼近，再估计外部的半范数尾，即可证明该交换合法。内层函数经普通 Fubini 化为
\[
 \int\varphi(x-y)\widehat\psi(x)\,\mathrm{d}x
 =\mathcal F(\widehat\varphi\,\psi)(y).
\]
因此 \(\langle\mathcal F(T*\varphi),\psi\rangle=\langle T,\mathcal F(\widehat\varphi\psi)\rangle=\langle\widehat T,\widehat\varphi\psi\rangle\)，即得结论。
\end{proof}

这里一般只得到多项式增长，不能断言卷积仍为 Schwartz 函数。例如 \(T=1\)、\(\int\varphi=1\) 时，\(T*\varphi=1\)。任意两个温和分布之间则没有普遍可定义的卷积；一方为 Schwartz 函数是本命题中保证配对及积分操作成立的条件。

% !TeX root = ../../Book2.tex
\section{常系数算子的基本解}
\label{b2:sec:0504}

常系数微分在频率侧变成多项式乘法。若能解出对点质量的响应，就可以通过卷积生成其他右端的解。这个响应称为基本解；它通常带有奇性，因此分布语言不可省略。

\begin{definition}\label{b2:def:05-fundamental}
设 \(P(\partial)=\sum_{|\alpha|\leq m}a_\alpha\partial^\alpha\) 为常系数微分算子。若分布 \(E\) 满足 \(P(\partial)E=\delta_0\)，则称 \(E\) 为该算子的\term[jiben-jie]{基本解}[fundamental solution]。
\end{definition}

\subsection{微分算子的符号}
\label{b2:subsec:050401}

定义其\term[weifen-suanzi-fuhao]{符号}[symbol]为
\[
 p(\xi)\coloneqq\sum_{|\alpha|\leq m}a_\alpha(2\pi\mathrm{i}\xi)^\alpha.
\]
由\eqref{b2:eq:05-distribution-operations}，\(\mathcal F(P(\partial)T)=p\widehat T\)。例如 \(-\Delta\) 的符号为 \(4\pi^2|\xi|^2\)，而 \(1-\Delta\) 的符号为 \(1+4\pi^2|\xi|^2\)。正号不同会直接改变能量与基本解，不能只记住“Laplacian 对应频率平方”。

\subsection{Fourier 方法}
\label{b2:subsec:050402}

若基本解温和，则其频率方程为 \(p\widehat E=1\)。当 \(p\) 处处非零，且 \(1/p\) 及其全部导数至多多项式增长时，可以定义 \(\widehat E=1/p\)，再取分布逆变换。这给出一个充分条件。若 \(p\) 有零点，就不能直接在这些点除法；即使几乎处处可以相除，也须检查所得函数是否局部可积。

\begin{proposition}\label{b2:prop:05-bessel-fundamental}
令
\[
 H_t(x)=(4\pi t)^{-d/2}e^{-|x|^2/(4t)},\quad
 E(x)=\int_0^\infty e^{-t}H_t(x)\,\mathrm{d}t.
\]
则 \(E\in L^1(\mathbb R^d)\)，非负且积分为 \(1\)，并在分布意义下满足 \((1-\Delta)E=\delta_0\)。
\end{proposition}
\begin{proof}
Gauss 积分给出 \(\int H_t=1\)，所以 Tonelli 说明 \(\int E=\int_0^\infty e^{-t}\,\mathrm{d}t=1\)，并保证定义积分几乎处处有限。由缩放公式，\(\widehat H_t(\xi)=e^{-4\pi^2t|\xi|^2}\)。再次按绝对积分 \(\int e^{-t}\|H_t\|_1\,\mathrm{d}t=1\) 换序，得到
\[
 \widehat E(\xi)=\int_0^\infty e^{-t(1+4\pi^2|\xi|^2)}\,\mathrm{d}t
 =\dfrac1{1+4\pi^2|\xi|^2}.
\]
因此 \(\mathcal F((1-\Delta)E)=1=\widehat\delta_0\)，分布变换的单射性给出结论。
\end{proof}

在一维还可直接写出 \(E(x)=\tfrac12e^{-|x|}\)：它在原点以外满足 \(E''=E\)，而一阶导数在原点的跳跃为 \(-1\)，所以 \(E''=E-\delta_0\)。这与频率分母 \(1+4\pi^2\xi^2\) 一致，也说明基本解方程的点质量来自导数的跳跃。

\subsection{Laplace 算子的基本解}
\label{b2:subsec:050403}

以下统一求 \(-\Delta\) 的基本解。记 \(\sigma_{d-1}\) 为单位球面的面积；这里的符号是表面积，不是单位球体积。

\begin{theorem}\label{b2:thm:05-laplace-fundamental}
局部可积函数
\[
 \Phi_d(x)=
 \begin{cases}
 \dfrac1{(d-2)\sigma_{d-1}}|x|^{2-d},&d\geq3,\\[4pt]
 -\dfrac1{2\pi}\log|x|,&d=2,\\[4pt]
 -\dfrac12|x|,&d=1
 \end{cases}
\]
定义温和分布，满足 \(-\Delta\Phi_d=\delta_0\)。
\end{theorem}
\begin{proof}
原点附近，\(d\geq3\) 时的绝对积分由 \(\int_0^1r^{2-d}r^{d-1}\,\mathrm{d}r=1/2\) 控制；二维由 \(\int_0^1r|\log r|\,\mathrm{d}r<\infty\) 控制。一维函数连续。无穷远只有幂或对数增长，故均满足温和性的加权积分条件。

当 \(d\geq2\) 时，在原点外用径向公式 \(\Delta v=v''+(d-1)v'/r\) 验证 \(\Delta\Phi_d=0\)。取 \(\varphi\in C_c^\infty\)，在去掉半径 \(\varepsilon\) 小球的区域上使用 Green 公式。外边界处测试函数及导数为零；内边界的区域外法向是 \(-e_r\)，故
\[
 \int_{|x|>\varepsilon}\Phi_d(-\Delta\varphi)\,\mathrm{d}x
 =\int_{|x|=\varepsilon}
 \left(-\varphi\,\partial_r\Phi_d+\Phi_d\partial_r\varphi\right)\,\mathrm{d}S.
\]
在两种维数情形中，都有 \(-\partial_r\Phi_d(\varepsilon)=1/(\sigma_{d-1}\varepsilon^{d-1})\)。第一项正好是 \(\varphi\) 的球面平均，趋于 \(\varphi(0)\)。第二项的绝对值至多为常数乘 \(\varepsilon\)（\(d\geq3\)）或 \(\varepsilon|\log\varepsilon|\)（\(d=2\)），故趋零。左侧由局部可积性趋于整个积分，得到分布恒等式。

一维时，分别在负、正半轴作两次分部积分，有 \(\int |x|\varphi''(x)\,\mathrm{d}x=2\varphi(0)\)，因此 \(-\Phi_1''=\delta_0\)。首先得到的等式在 \(\mathcal D\) 上成立；两边均温和，利用 \(\mathcal D\) 在 \(\mathcal S\) 中稠密便推广到 \(\mathcal S'\)。
\end{proof}

二维的对数和一维的线性增长，正好提醒我们：\(1/(4\pi^2|\xi|^2)\) 在这些维数的原点并不可积，不能用该函数的普通积分冒充基本解的频率分布。上面的空间构造已经定义了合法的温和基本解，其变换满足 \(4\pi^2|\xi|^2\widehat\Phi_d=1\)；有奇性的频率乘法方程应在这个意义下理解。

\subsection{基本解与方程求解}
\label{b2:subsec:050404}

\begin{proposition}\label{b2:prop:05-fundamental-convolution}
设温和分布 \(E\) 满足 \(P(\partial)E=\delta_0\)，且 \(f\in C_c^\infty(\mathbb R^d)\)。则
\[
 u(x)\coloneqq\langle E_y,f(x-y)\rangle
\]
是光滑函数，并满足 \(P(\partial)u=f\)。
\end{proposition}
\begin{proof}
对固定 \(x\)，函数 \(y\mapsto f(x-y)\) 属于 \(\mathcal S\)。当 \(x\) 限于紧集时，这族函数有共同紧支集，并且关于 \(x\) 的差商在每个 Schwartz 半范数中趋于相应导数；这可对各阶导数使用 Taylor 公式的一致余项估计。连续泛函 \(E\) 可以通过这些极限，所以 \(u\) 光滑，导数可移到 \(f\) 上。

又因 \(\partial_x^\alpha f(x-y)=(-1)^{|\alpha|}\partial_y^\alpha f(x-y)\)，分布导数定义给出
\[
 P(\partial)u(x)=\langle P(\partial)E_y,f(x-y)\rangle
 =\langle\delta_0,f(x-\cdot)\rangle=f(x).
\]
\end{proof}

特别地，\(u=\Phi_d*f\) 解 \(-\Delta u=f\)。这只解决全空间的方程，没有自动满足某个有界区域上的 Dirichlet 或 Neumann 条件。基本解也一般不唯一：加上任何齐次方程的温和解仍为基本解。例如给 \(\Phi_d\) 加一个调和多项式不改变 \(-\Delta\Phi_d=\delta_0\)。要获得唯一的边值解，还须结合区域、边界条件与增长约束。



\ProseParagraphBreak

二维基本解还表现出临界尺度的差别。对 \(\lambda>0\)，
\[
 \Phi_2(\lambda x)=\Phi_2(x)-\dfrac{\log\lambda}{2\pi}.
\]
它不严格齐次，而是相差一个常数；该常数的 Laplace 为零，所以仍与基本解方程相容。高维的 \(\Phi_d\) 则按次数 \(2-d\) 齐次。这说明从一般维数公式形式代入 \(d=2\) 不能直接得到答案：分母与幂次同时退化，要保留对数结构及可加的齐次解。


\begin{chaptersummary}
Schwartz 空间用所有阶导数和所有多项式权重控制函数。Fourier 变换在其中交换微分与乘法，反演使它成为连续自同构；Hermite 基则把这一变换的四种特征值具体呈现出来。温和分布通过作用于测试函数定义，双线性配对把已有的变换公式连续延伸到 Dirac 测度、多项式以及全部 \(L^p\) 函数。

\ProseParagraphBreak

常系数方程的频率表达是多项式乘法，但零点处不能随意相除。基本解的构造必须兼顾局部可积性、原点质量和无穷远增长。负 Laplace 的幂核、对数核和一维线性核各反映一种维数边界；与紧支集光滑右端卷积可解全空间方程，却不自动决定边值条件。下一章将用同一测试空间把离散格点与连续频率联系起来。
\end{chaptersummary}

\BookExercises

以下题目把半范数、差商和频率符号用于具体对象。前几题强调收敛与符号，后几题处理奇性、频率特征结构和方程。

% 自拟：所有导数一致趋零仍不控制多项式权重。
\begin{exercise}\label{b2:ex:05-schwartz-convergence}
设非零 \(\rho\in C_c^\infty(\mathbb R)\)，令 \(\varphi_n(x)=n^{-1}\rho(x/n)\)。证明每个固定阶导数都一致趋于零，但 \(\varphi_n\) 不在 \(\mathcal S\) 中趋零。它作为温和分布是否趋零？
\end{exercise}
\begin{solution}
有 \(\|\varphi_n^{(k)}\|_\infty=n^{-k-1}\|\rho^{(k)}\|_\infty\to0\)。选 \(x_0\ne0\) 使 \(\rho(x_0)\ne0\)，则 \(p_1(\varphi_n)\geq n^{-1}(1+n|x_0|)|\rho(x_0)|\)，下极限为正，故不按 Schwartz 半范数趋零。另一方面，对 \(\psi\in\mathcal S\)，\(|\int\varphi_n\psi|\leq n^{-1}\|\rho\|_\infty\cdot\|\psi\|_1\to0\)，所以在 \(\mathcal S'\) 中趋零。
\end{solution}

% 自拟：差分点质量的极限与傅里叶符号同时核对。
\begin{exercise}\label{b2:ex:05-delta-differences}
在一维求下列分布列的极限，并计算极限的 Fourier 变换：
\[
 n(\delta_{1/n}-\delta_0),\quad
 n^2(\delta_{1/n}+\delta_{-1/n}-2\delta_0).
\]
\end{exercise}
\begin{solution}
对 \(\varphi\in\mathcal S\) 使用一阶与二阶 Taylor 差商，配对极限分别为 \(\varphi'(0)\) 和 \(\varphi''(0)\)，故极限为 \(-\delta_0'\)、\(\delta_0''\)。它们的变换分别为 \(-2\pi\mathrm{i}\xi\)、\(-4\pi^2\xi^2\)。也可从指数差商核对：\(n(e^{-2\pi\mathrm{i}\xi/n}-1)\to-2\pi\mathrm{i}\xi\)，其模由 \(2\pi|\xi|\) 控制；对称二阶差商的模由 \(4\pi^2\xi^2\) 控制，故两者都可对 Schwartz 测试函数用控制收敛。
\end{solution}

% 自拟：乘以x的核及原点导数项，训练分布除法。
\begin{exercise}\label{b2:ex:05-x-times-distribution}
在 \(\mathcal S'(\mathbb R)\) 中证明 \(xT=0\) 当且仅当 \(T=c\delta_0\)。再求方程 \(xT=\delta_0\) 的全部解。
\end{exercise}
\begin{solution}
取 \(\rho\in C_c^\infty\)，在原点附近为 \(1\)。任意 \(\varphi\in\mathcal S\) 可写成 \(\varphi=\varphi(0)\rho+x\psi\)，其中 \(\psi=(\varphi-\varphi(0)\rho)/x\) 在原点用光滑延拓，远处除以 \(x\) 保持快速下降，所以 \(\psi\in\mathcal S\)。若 \(xT=0\)，便有 \(\langle T,\varphi\rangle=\varphi(0)\langle T,\rho\rangle\)，即 \(T=c\delta_0\)。反向直接成立。又因 \(\langle x\delta_0',\varphi\rangle=-(x\varphi)'(0)=-\varphi(0)\)，一个特解为 \(-\delta_0'\)。全部解因而为 \(-\delta_0'+c\delta_0\)。
\end{solution}

% 自拟：主值分布、奇偶性和一阶分布方程组合得到变换常数。
\begin{optionalexercise}\label{b2:ex:05-principal-value}
证明
\[
 \langle T,\varphi\rangle
 =\lim_{\varepsilon\downarrow0}\int_{|x|>\varepsilon}\dfrac{\varphi(x)}x\,\mathrm{d}x
\]
定义温和分布，称为 \(\operatorname{pv}(1/x)\)。证明 \(xT=1\)，并计算 \(\widehat T\)。
\end{optionalexercise}
\begin{solution}
在 \(|x|<1\) 内减去 \(\varphi(0)\)，常数项因对称性抵消，剩余绝对值由 \(\|\varphi'\|_\infty\) 控制；在 \(|x|\geq1\) 内用 \(p_2(\varphi)(1+|x|)^{-2}/|x|\) 控制。因此极限存在且满足有限半范数界。代入 \(x\varphi\) 得 \(xT=1\)，并可换元看出 \(T\) 为奇分布。

由变换公式，\(\frac{\mathrm{i}}{2\pi}\partial_\xi\widehat T=\delta_0\)，故 \(\partial_\xi\widehat T=-2\pi\mathrm{i}\delta_0\)。函数 \(-\pi\mathrm{i}\operatorname{sgn}\xi\) 的分布导数正好为右端，所以两者之差的导数为零。导数为零的温和分布是常数：对积分为零的 \(\psi\in\mathcal S\)，其原函数 \(\int_{-\infty}^x\psi\) 属于 \(\mathcal S\)，故该分布在所有零积分测试函数上为零；再减去一个固定单位积分测试函数的倍数即可。于是 \(\widehat T=-\pi\mathrm{i}\operatorname{sgn}\xi+c\)。Fourier 变换保持奇偶性，而 \(T\) 为奇，所以 \(c=0\)。
\end{solution}

% 自拟：求Hermite低阶式，并把频率特征关系与谐振子联系。
\begin{exercise}\label{b2:ex:05-hermite-oscillator}
求 \(h_1,h_2\) 的显式公式，证明
\[
 \left(-\dfrac{\mathrm{d}^2}{\mathrm{d}x^2}+4\pi^2x^2\right)h_n
 =(4n+2)\pi h_n,
\]
并把 \(x^2h_0\) 展成 Hermite 函数的有限线性组合。
\end{exercise}
\begin{solution}
从上升算子计算得 \(h_1=\sqrt{4\pi}\,xh_0\)、\(h_2=2^{-1/2}(4\pi x^2-1)h_0\)。展开算子有 \(4\pi a^*a=-\mathrm{d}^2/\mathrm{d}x^2+4\pi^2x^2-2\pi\)。由 \(a^*ah_n=nh_n\) 得到特征方程。第二个显式公式给出 \(x^2h_0=(h_0+\sqrt2h_2)/(4\pi)\)。其变换为 \((h_0-\sqrt2h_2)/(4\pi)\)，也可与两次频率求导公式核对。
\end{solution}

% 自拟：有界光滑乘子若导数失控，可以破坏S。
\begin{exercise}\label{b2:ex:05-fast-multiplier}
在一维令 \(a(x)=e^{\mathrm{i}e^{x^2}}\)、\(\varphi(x)=e^{-x^2}\)。证明 \(a\) 光滑且有界，\(\varphi\in\mathcal S\)，但 \(a\varphi\notin\mathcal S\)。比较乘子 \(b(x)=e^{\mathrm{i}x^3}\) 的行为。
\end{exercise}
\begin{solution}
直接求导得 \((a\varphi)'=a(-2xe^{-x^2}+2\mathrm{i}x)\)，其模随 \(|x|\) 无界，故不满足 Schwartz 条件。对 \(b\)，每个导数都是一个多项式乘 \(b\)，因而至多多项式增长。由正文的乘法命题，乘以 \(b\) 连续保持 \(\mathcal S\)；其逆乘子 \(\overline b\) 也有相同性质，所以还给出连续自同构。
\end{solution}

% 自拟：尺度决定屏蔽算子的基本解与L1质量。
\begin{exercise}\label{b2:ex:05-resolvent-scale}
设 \(E\) 为正文构造的 \(1-\Delta\) 基本解。对 \(a>0\)，证明 \(E_a(x)=a^{d/2-1}E(\sqrt a\,x)\) 是 \(a-\Delta\) 的可积基本解，求其变换和 \(L^1\) 范数。若 \(f\in L^p\)，证明 \(\|E_a*f\|_p\leq a^{-1}\|f\|_p\)。
\end{exercise}
\begin{solution}
尺度公式给出
\[
 \widehat E_a(\xi)=a^{-1}\widehat E(\xi/\sqrt a)
 =\dfrac1{a+4\pi^2|\xi|^2}.
\]
乘以符号 \(a+4\pi^2|\xi|^2\) 后得到 \(1\)，故为基本解。它非负，换元得 \(\|E_a\|_1=a^{-1}\|E\|_1=a^{-1}\)。最后用 Young 卷积不等式。这个估计允许任意 \(1\leq p\leq\infty\)，但不单凭此范数界宣称解具有任意阶强导数。
\end{solution}

% 自拟：离散测度的温和性及其不可微函数式解释。
\begin{exercise}\label{b2:ex:05-polynomial-comb}
设 \((c_k)_{k\in\mathbb Z^d}\) 满足 \(|c_k|\leq C(1+|k|)^M\)。证明 \(T=\sum_{k\in\mathbb Z^d}c_k\delta_k\) 由绝对收敛配对定义温和分布，且立方体内的有限截断在 \(\mathcal S'\) 中趋向 \(T\)。
\end{exercise}
\begin{solution}
选整数 \(N>M+d\)。则
\[
 \sum_k|c_k\varphi(k)|
 \leq C p_N(\varphi)\sum_k(1+|k|)^{M-N}<\infty.
\]
最后的格点级数可按整数半径或二进环带比较体积：半径约 \(2^j\) 的点数至多为 \(C_d2^{jd}\)，所以所得几何级数收敛。这同时给出有限半范数界。对每个固定测试函数，绝对收敛级数的截断尾部趋零，正好是 \(\mathcal S'\) 收敛。
\end{solution}
